\chapter{Introducción}

Este capítulo sitúa el problema de investigación, delimita una especificación verificable para el prototipo y adelanta las contribuciones, el alcance y la organización de la memoria. La intención es distinguir desde el inicio entre la evidencia que puede obtenerse mediante experimentación fuera de línea (\emph{offline}) y cualquier afirmación de preparación para un despliegue operativo.

\section{Motivación}
Las infraestructuras de red modernas generan volúmenes de tráfico que superan con creces lo que los analistas humanos pueden inspeccionar manualmente. Una única red empresarial puede producir millones de registros de flujo al día, y los despliegues a gran escala agregan flujos procedentes de miles de endpoints, servicios en la nube y enlaces entre centros de datos \parencite{ScarfoneMell2007}. Distinguir la comunicación legítima de la actividad maliciosa dentro de este volumen es uno de los desafíos operativos centrales en la seguridad de redes. Las respuestas tradicionales se apoyan en sistemas de detección de intrusiones en red (NIDS, del inglés \emph{Network Intrusion Detection Systems}) que aplican coincidencia de firmas, umbrales estadísticos o reglas específicas de protocolo para señalar actividad sospechosa. Estos enfoques son efectivos frente a patrones de amenaza conocidos, pero requieren actualizaciones manuales constantes de reglas, resultan en gran medida ineficaces frente a variantes novedosas de ataque y generan volúmenes de alertas que habitualmente desbordan a los equipos de seguridad.

El aprendizaje automático (ML, del inglés \emph{Machine Learning}) ofrece una perspectiva complementaria. En lugar de codificar el conocimiento experto como reglas estáticas, los modelos de ML pueden aprender patrones estadísticos a partir de datos de tráfico etiquetados y generalizar hacia instancias de ataque no vistas previamente. A medida que el campo ha madurado, las representaciones a nivel de flujo se han convertido en un sustrato práctico para este aprendizaje: cada conexión de red se resume como un conjunto fijo de características estadísticas que cubren duración, conteos de paquetes, totales de bytes, tiempos entre llegadas, distribuciones de flags TCP y métricas basadas en tasas \parencite{Sperotto2010FlowIDS,Umer2017FlowIDS}. Esta reducción hace abordable el aprendizaje a gran escala y evita las complicaciones legales y técnicas de la inspección profunda de paquetes, que resulta cada vez más ineficaz de todos modos a medida que crece la fracción de tráfico cifrado en internet.

La limitación clave del aprendizaje supervisado estándar en este contexto es que trata la detección de intrusiones como un problema de clasificación simétrico, minimizando el error de clasificación global sin tener en cuenta las consecuencias operativas de los distintos tipos de error \parencite{Axelsson1999BaseRate}. En la práctica, esas consecuencias son profundamente asimétricas. Un falso negativo permite a un actor de amenazas avanzar sin ser detectado y puede conducir a exfiltración de datos, despliegue de ransomware o interrupción del servicio. Un falso positivo genera una alerta o un bloqueo de tráfico que consume tiempo del analista, pero rara vez causa daños duraderos. El aprendizaje por refuerzo (RL, del inglés \emph{Reinforcement Learning}) proporciona un marco natural para codificar esta asimetría directamente en el objetivo de aprendizaje: al diseñar una función de recompensa que penaliza los ataques no detectados más severamente que las falsas alarmas, un agente de RL puede ser guiado hacia políticas que reflejen las prioridades de riesgo operativo real en lugar de la exactitud agregada \parencite{SuttonBarto2018RL,Lee2002CostSensitiveIDS}. Explorar cómo un agente basado en la \emph{Quantile Regression Deep Q-Network} (QRDQN) navega estos compromisos sobre datos a nivel de flujo es la motivación central de esta tesis \parencite{Dabney2018QRDQN}.

\section{Definición del problema}

\subsection{Desafíos metodológicos}
Aplicar el aprendizaje por refuerzo a la detección de intrusiones en red plantea varios desafíos metodológicos que van más allá de sustituir un algoritmo de RL por un clasificador estándar. El primer desafío es la calidad de los datos. Los bancos de pruebas públicos (\emph{benchmarks}) de NIDS como CICIDS2017 son ampliamente utilizados, pero contienen artefactos conocidos de su flujo de procesamiento (\emph{pipeline}) de extracción de flujos, entre ellos ruido en las etiquetas, características duplicadas, fallos de implementación de CICFlowMeter y campos estadísticos que codifican inadvertidamente el contexto del día de captura \parencite{Ring2019DatasetSurvey,Engelen2021CICIDSIssues,Lanvin2023Errors}. Los modelos entrenados sin auditar estos artefactos pueden aprender a reconocer firmas de extracción o identificadores específicos de escenario en lugar de comportamiento de ataque genuino, produciendo puntuaciones infladas en el \emph{benchmark} que no se transfieren a ningún entorno nuevo.

El segundo desafío es la fuga de información. Muchos estudios publicados de NIDS evalúan los modelos utilizando particiones aleatorias por filas sobre todo el tráfico disponible, lo que puede situar flujos altamente correlacionados a ambos lados del límite entrenamiento-prueba. Más allá de las filas duplicadas, la fuga puede introducirse a través de direcciones IP, marcas temporales absolutas, identificadores del flujo \textit{(Flow IDs)} y números de puerto que son específicos de un escenario de ataque programado en lugar de indicativos del comportamiento malicioso en general \parencite{Arp2020DosDontsMLSecurity,Kaufman2011LeakageDataMining}. Los pasos de preprocesamiento como el escalado de características, aplicados globalmente antes de la partición, pueden contaminar adicionalmente la evaluación al permitir que la información distribucional de la partición de prueba influya en la normalización del entrenamiento. Un estudio riguroso no puede limitarse a elegir un algoritmo de alto rendimiento y reportar su exactitud sobre un conjunto de datos público dividido aleatoriamente.

El tercer desafío es la comparación justa contra una línea base (\emph{baseline}) de un método estándar. Los métodos de aprendizaje por refuerzo profundo (DRL, del inglés \emph{Deep Reinforcement Learning}) pueden ser costosos de ajustar, sensibles a los hiperparámetros y más lentos en converger que los modelos supervisados tabulares. El bosque aleatorio (\emph{Random Forest}) y los ensamblados de árboles con potenciación de gradiente (\emph{gradient boosting}) son líneas base sólidas para la detección de intrusiones a nivel de flujo y pueden alcanzar puntuaciones casi perfectas en CICIDS2017 bajo condiciones de evaluación favorables \parencite{LiuLang2019IDSSurvey,Apruzzese2022CrossEvaluationNIDS}. Un defensor basado en RL no puede evaluarse de manera significativa sin una comparación bajo el mismo protocolo frente a dichas líneas base. La pregunta no es si QRDQN alcanza alta exactitud en una sola ejecución, sino si su formulación sensible al coste ofrece alguna ventaja empírica sobre un modelo supervisado bien ajustado bajo esquemas de características, protocolos de partición y métricas de evaluación equivalentes.

\subsection{Especificación verificable del proyecto}
\label{sec:especificacion-proyecto}
Como adaptación al contexto de investigación de la especificación propia de un proyecto de desarrollo de software (PDS), esta tesis formula el trabajo como un sistema experimental verificable y no como un producto de bloqueo listo para operar. La Tabla~\ref{tab:especificacion-proyecto} fija el problema que se aborda, los límites que preservan la seguridad metodológica y la evidencia exigida para considerar cubierto cada compromiso.

\begin{table}[htbp]
\centering
\small
\begin{tabular}{@{}p{2.6cm}p{9.7cm}@{}}
\toprule
\multicolumn{2}{@{}p{12.3cm}@{}}{\textbf{Especificación verificable del prototipo experimental}} \\
\midrule
Necesidad & Construir y evaluar un defensor offline que clasifique flujos de red como \texttt{PERMIT} o \texttt{BLOCK}, haciendo explícita la asimetría entre permitir un ataque y bloquear tráfico legítimo. \\
Entradas & CSV de flujos de CICIDS2017 y, para la comprobación complementaria, flujos extraídos de tráfico de laboratorio; todos se adaptan al contrato canónico y se procesan sin usar identificadores ni otros campos propensos a fuga. \\
Proceso & Limpieza y escalado ajustados solo con entrenamiento, representación de $76+76=152$ dimensiones, entrenamiento QRDQN, comparación con Random Forest y escalera de validación con artefactos persistidos. \\
Salidas & Predicciones binarias experimentales, métricas por clase y artefactos de ejecución trazables; \texttt{BLOCK} permanece como salida de inferencia offline y no acciona un cortafuegos ni una red real. \\
Verificación & Los compromisos C1--C4, especificados en las Tablas~\ref{tab:obj-c1} a~\ref{tab:obj-c4}, requieren código auditable, ejecuciones comprometidas y resultados interpretados dentro de las restricciones declaradas. \\
Límite de la afirmación & El rendimiento sobre CICIDS2017 y sobre una captura concreta de laboratorio constituye evidencia experimental acotada; no certifica generalización a producción ni bloqueo activo seguro. \\
\bottomrule
\end{tabular}
\caption[Especificación verificable del proyecto]{Especificación verificable del prototipo experimental, adaptada al contexto de un proyecto de desarrollo de software. Fuente: elaboración propia.}
\label{tab:especificacion-proyecto}
\end{table}

\section{Enfoque propuesto}
Esta tesis formula la tarea defensiva como un problema de aprendizaje por refuerzo offline con el conjunto de datos como entorno \parencite{Levine2020OfflineRL}. La elección del entorno offline es deliberada: desplegar un agente de RL que deba interactuar con una red en vivo para aprender generaría riesgos de seguridad inaceptables, ya que un agente en exploración tomaría inevitablemente acciones subóptimas durante el entrenamiento. Al tratar un conjunto de datos etiquetado estático como el entorno, el agente puede entrenarse y evaluarse sin ninguna exposición a una red en vivo. El tráfico de red se procesa en un esquema canónico de flujo de 76 características fijo, independiente del formato nativo de cualquier conjunto de datos específico. Este contrato canónico separa la nomenclatura de columnas propia de la extracción del espacio de observación del modelo y hace posible procesar tanto flujos de CICIDS2017 como tráfico capturado en un laboratorio doméstico cerrado a través del mismo pipeline sin modificación alguna.

Para gestionar esquemas heterogéneos al transitar entre herramientas de extracción, se añade una máscara de presencia/ausencia de 76 valores, produciendo una observación de 152 dimensiones. Cada posición indica si existe una columna fuente mapeada para la característica canónica correspondiente. En CICIDS2017 las 76 columnas canónicas están cubiertas y la máscara es constante; en la Fase~2 puede revelar columnas no disponibles. Los valores no finitos de una columna existente se sanean antes del mapeo y no cambian su indicador. Esta semántica, más estrecha que una máscara de validez fila a fila, coincide con la implementación actual y hace auditable la compatibilidad estructural entre extractores.

La tarea de decisión se mapea a un espacio de acciones binario: \texttt{PERMIT} para el tráfico clasificado como benigno, y \texttt{BLOCK} para el tráfico clasificado como ataque. Un agente \textit{Quantile Regression Deep Q-Network} \parencite{Dabney2018QRDQN}, implementado a través de la biblioteca \texttt{sb3-contrib} \parencite{SB3ContribQRDQNDocs}, se entrena dentro de un entorno compatible con Gymnasium \parencite{GymnasiumDocs}. El entorno aplica una función de recompensa asimétrica que penaliza los falsos negativos más severamente que los falsos positivos, codificando directamente las prioridades de riesgo operativo descritas anteriormente \parencite{Cardenas2006IDSFramework}. El conjunto de datos CICIDS2017 \parencite{Sharafaldin2018CICIDS2017,CICIDS2017Web} sirve como benchmark interno primario, procesado a través de un pipeline de carga curado que excluye los campos propensos a fuga de información antes de que comience cualquier entrenamiento del modelo. El rendimiento se evalúa frente a una línea base de Random Forest bajo el mismo esquema de características y protocolo de partición \parencite{Engelen2021CICIDSIssues}. Un segundo nivel de validación extiende la evaluación ejecutando inferencia offline sobre características de flujo extraídas de tráfico capturado por el operador en un laboratorio doméstico cerrado y no adversarial, ofreciendo una indicación parcial y de validez externa limitada del comportamiento del sistema ante un cambio de dominio.

\section{Contribución de la tesis}
Las contribuciones se expresan mediante los cuatro compromisos contractuales del Capítulo~2. Esta formulación permite separar el producto construido, la evidencia de evaluación y los límites de cada conclusión, en lugar de presentar el rendimiento de una sola partición como una demostración global.

\textbf{C1 --- Pipeline experimental reproducible con contrato de datos canónico.} La primera contribución es un pipeline offline de extremo a extremo que fija una representación de $76+76=152$ dimensiones y conserva los artefactos necesarios para reconstruir cada resultado. La Tabla~\ref{tab:obj-c1} especifica este compromiso y el Anexo~\ref{cap:anexo-artefactos} documenta el contrato de evidencia.

\textbf{C2 --- Defensor RL y jerarquía de validación con control de fugas.} La segunda contribución es la implementación de QRDQN como defensor binario y el diseño de una campaña que separa MAIN aleatoria, día completo, tamaño, semilla y cuatro escenarios retenidos, junto con controles directos, barajado de etiquetas y duplicados. La Tabla~\ref{tab:obj-c2} fija el compromiso; el Capítulo~\ref{cap:protocolo-experimental} distingue el diseño aprobado de la evidencia todavía pendiente.

\textbf{C3 --- Decisión sensible al coste con lectura multiobjetivo.} La tercera contribución convierte la asimetría entre ataques permitidos y tráfico legítimo bloqueado en un elemento auditable del diseño, no en una interpretación posterior de la exactitud. La Tabla~\ref{tab:obj-c3} enlaza esta contribución con la matriz de recompensa y su lectura de seguridad e impacto operativo (Tabla~\ref{tab:recompensa-multiobjetivo}).

\textbf{C4 --- Comparación honesta y transferencia offline acotada.} La cuarta contribución incorpora Random Forest bajo particiones emparejadas y una Fase~2 que reutiliza modelo y preprocesamiento sin reajuste. La evidencia histórica se conserva como antecedente; la comparación final y la inferencia fresca dependen de la campaña. La Tabla~\ref{tab:obj-c4} fija el compromiso y la Tabla~\ref{tab:estructura-resultados-finales} define la evidencia necesaria, sin atribuir a \texttt{BLOCK} ningún efecto operativo.

\section{Alcance y limitaciones}
El alcance de esta investigación se limita estrictamente a la evaluación offline, tabular y a nivel de flujo. Las salidas \texttt{PERMIT} y \texttt{BLOCK} son clasificaciones binarias experimentales sobre registros de flujo estáticos; el sistema no realiza bloqueo en vivo de cortafuegos, descarte de paquetes, manipulación de \texttt{iptables} ni prevención de intrusiones en línea. El entorno no simula un adversario activo cuyas acciones posteriores dependan de las decisiones del defensor, lo que significa que el sistema no opera como un plano de control de toma de decisiones secuenciales completo. Estas restricciones son decisiones de diseño deliberadas que hacen que la evaluación sea abordable, reproducible y segura. No son omisiones que deban subsanarse en iteraciones futuras de este mismo prototipo.

CICIDS2017 se utiliza como benchmark interno primario para la experimentación controlada. Los resultados sobre este conjunto de datos deben interpretarse como evidencia sobre el comportamiento bajo un benchmark específico y bien caracterizado, no como prueba de generalización a redes de producción arbitrarias. El tráfico de laboratorio, generado por el operador en un entorno doméstico cerrado y no adversarial, se utiliza como nivel de evaluación secundario para examinar el cambio de dominio; dada su procedencia, su validez externa es limitada y, además, también es inferencia offline y no despliegue en vivo. La tesis ocupa una posición metodológica entre la evaluación pura sobre benchmark y la validación de despliegue completo, apuntando a la reproducibilidad e interpretabilidad del marco experimental más que a la operatividad en producción.

\section{Estructura del documento}
Tras esta introducción, el Capítulo~2 fija los objetivos contractuales y específicos, el alcance, las restricciones y la planificación. El Capítulo~3 revisa el estado del arte de los NIDS basados en flujos, los conjuntos de datos públicos, el aprendizaje automático y el aprendizaje por refuerzo aplicados a detección de intrusiones, junto con sus riesgos metodológicos. El Capítulo~4 presenta el diseño del sistema: datos, limpieza anti-fuga, contrato canónico, formulación del entorno y agente QRDQN. El Capítulo~5 define el protocolo experimental, las particiones, las métricas, la línea base y la inferencia offline de la Fase~2.

El Capítulo~6 fija la estructura de incorporación de la campaña final y conserva, claramente separada, la evidencia histórica previa. El Capítulo~7 organiza la discusión por generalización, duplicados, coste, comparación de modelos, escala, semillas, escenarios y cambio de dominio; el Capítulo~8 expone las limitaciones y el Capítulo~9 delimita las implicaciones éticas. El Capítulo~10 presenta una conclusión provisional que podrá cerrarse desde los artefactos finales. Los anexos reúnen reproducción, entorno, esquema canónico y catálogo de artefactos.
